\documentclass[11pt]{article}
\usepackage[margin=1in]{geometry}
\usepackage{amsmath,amssymb,amsthm}
\usepackage{graphicx}
\usepackage{microtype}
\usepackage[colorlinks=true,linkcolor=blue,citecolor=blue,urlcolor=blue]{hyperref}
\usepackage{enumitem}

\newtheorem{theorem}{Theorem}
\newtheorem{lemma}[theorem]{Lemma}
\newcommand{\K}{\mathcal K}
\newcommand{\Lop}{\mathcal L}
\newcommand{\Pw}{\mathcal P_w}
\newcommand{\Poly}{\mathcal P}

\title{Critical-Degree \(M\)-Ideals for Finite-Dimensional Ranges\\
\large A partial answer to arXiv:1102.3850}
\author{Run \texttt{fa\_banach\_001}, lane 17\\Model GPT5.6}
\date{11 August 2026}

\begin{document}
\maketitle

\noindent\textbf{Status.}
Candidate partial result, likely valid. This packet solves the equality
endpoint for every nonzero finite-dimensional range space \(F\); it does not
settle the source question for arbitrary property-\((M)\) spaces \(F\).

\section{Source question}

The source is V. Dimant and S. Lassalle,
\emph{\(M\)-structures in vector-valued polynomial spaces}, published in
Journal of Convex Analysis 19 (2012), 685--711 \cite{DL}. Its arXiv id is
1102.3850.
On PDF page 22, immediately before Corollary 4.9, the authors state:

\begin{quote}
\emph{We do not know if the result remains true even for the case
\(n=cd(E,F)=cd(E)\).}
\end{quote}

Here the result under discussion is the following lifting statement. If
\(\K(E)\) is an \(M\)-ideal in \(\Lop(E)\), \(F\) has property \((M)\), and
\[
 n=cd(E,F)<cd(E),
\]
then \(\Pw(^nE,F)\) is an \(M\)-ideal in \(\Poly(^nE,F)\).

\begin{center}
\includegraphics[width=0.96\textwidth]{figures/open_problem_crop.png}
\end{center}

\section{Definitions and imported results}

The critical degree \(cd(E,F)\), when it exists, is the least positive
integer \(k\) for which some continuous \(F\)-valued \(k\)-homogeneous
polynomial on \(E\) is not weakly continuous on bounded sets. We write
\(cd(E)=cd(E,\mathbb K)\).

A pair \((E,F)\) has the \(n\)-polynomial property \((M)\) when every
\(P\in\Poly(^nE,F)\) with \(\|P\|\le1\) satisfies
\[
 \limsup_\alpha\|v+P(x_\alpha)\|
 \le
 \limsup_\alpha\|u+x_\alpha\|^n
 \tag{1}\label{eq:vectorM}
\]
whenever \(\|v\|\le\|u\|^n\) and \((x_\alpha)\) is a bounded weakly
null net in \(E\).

We use three results already proved in the cited papers.

\begin{enumerate}[label=(\roman*)]
\item If \(\K(E)\) is an \(M\)-ideal in \(\Lop(E)\), then \(E\) has
property \((M)\) and admits the compact approximation net appearing in
Theorem 4.7 of \cite{DL}.
\item If \(n=cd(E)\) and \(E\) has property \((M)\), then \(E\) has the
scalar \(n\)-polynomial property \((M)\)
\cite[Proposition 3.10]{D}.
\item At \(n=cd(E,F)\), the compact approximation hypotheses in (i) and
the vector \(n\)-polynomial property \((M)\) imply that
\(\Pw(^nE,F)\) is an \(M\)-ideal in \(\Poly(^nE,F)\)
\cite[Theorem 4.7]{DL}.
\end{enumerate}

\section{Proof intuition}

The only new issue is lifting a scalar property-\((M)\) inequality to a
vector norm. In finite dimensions a norm can be scalarized coherently along
the net: choose a norming functional at each point, then pass to a subnet on
which those functionals converge in norm. The moving scalarizations can
therefore be replaced by one fixed functional \(f\), reducing the desired
inequality to scalar property \((M)\) for \(f\circ P\).

This mechanism is unavailable in general infinite-dimensional ranges:
weak-star convergence of norming functionals does not control their diagonal
pairing with the moving vectors \(P(x_\alpha)\).

\section{Formal proof}

\begin{lemma}\label{lem:cd}
If \(F\ne\{0\}\) is finite-dimensional and \(cd(E)\) exists, then
\[
 cd(E,F)=cd(E).
\]
\end{lemma}

\begin{proof}
Fix \(0\ne y\in F\). If a scalar \(k\)-homogeneous polynomial \(p\) is
not weakly continuous on bounded sets, neither is \(x\mapsto p(x)y\).
Thus \(cd(E,F)\le cd(E)\).

Conversely, let \(k<cd(E)\) and \(P\in\Poly(^kE,F)\). For every
\(f\in F^*\), the scalar polynomial \(f\circ P\) is weakly continuous
on bounded sets. Choose a finite basis \(f_1,\dots,f_m\) of \(F^*\).
If \(x_\alpha\to x\) weakly in a bounded subset of \(E\), then
\(f_j(P(x_\alpha)-P(x))\to0\) for each \(j\). Coordinate convergence
in finite-dimensional \(F\) is norm convergence, so
\(P(x_\alpha)\to P(x)\) in norm. Hence every such \(P\) is weakly
continuous on bounded sets and \(cd(E,F)\ge cd(E)\).
\end{proof}

\begin{lemma}\label{lem:lift}
Let \(F\) be finite-dimensional. If \(E\) has scalar
\(n\)-polynomial property \((M)\), then \((E,F)\) has vector
\(n\)-polynomial property \((M)\).
\end{lemma}

\begin{proof}
Let \(P\in\Poly(^nE,F)\) have norm at most one. Fix \(u\in E\),
\(v\in F\) with \(\|v\|\le\|u\|^n\), and a bounded weakly null net
\((x_\alpha)\). Put
\[
 A=\limsup_\alpha\|v+P(x_\alpha)\|,\qquad
 B=\limsup_\alpha\|u+x_\alpha\|^n.
\]
Pass to a subnet, still denoted \((x_\alpha)\), on which the first
norm converges to \(A\). For every \(\alpha\), Hahn--Banach gives
\(f_\alpha\in S_{F^*}\) such that
\[
 f_\alpha(v+P(x_\alpha))=\|v+P(x_\alpha)\|.
 \tag{2}\label{eq:norming}
\]
This choice is valid over both real and complex scalars.

The unit sphere of \(F^*\) is norm compact. After passing to a further
subnet, \(f_\alpha\to f\) in norm for some \(f\in B_{F^*}\). The net
\((v+P(x_\alpha))\) is bounded, so \eqref{eq:norming} implies
\[
 |f(v)+(f\circ P)(x_\alpha)|\longrightarrow A.
\]
The scalar polynomial \(f\circ P\) has norm at most one and
\[
 |f(v)|\le\|v\|\le\|u\|^n.
\]
Scalar \(n\)-polynomial property \((M)\) therefore yields
\[
 A
 \le \limsup_\alpha\|u+x_\alpha\|^n
 \le B.
\]
The last inequality holds because the selected net is a subnet of the
original net. This is exactly \eqref{eq:vectorM}.
\end{proof}

\begin{theorem}[Finite-dimensional endpoint]\label{thm:main}
Let \(E\) be a Banach space whose critical degree \(n=cd(E)\) exists,
and let \(F\ne\{0\}\) be finite-dimensional. If \(\K(E)\) is an
\(M\)-ideal in \(\Lop(E)\), then
\[
 \Pw(^nE,F)\quad\hbox{is an \(M\)-ideal in}\quad\Poly(^nE,F).
\]
In particular, Corollary 4.9 of \cite{DL} remains true at
\(n=cd(E,F)=cd(E)\) for every finite-dimensional \(F\).
\end{theorem}

\begin{proof}
The \(M\)-ideal hypothesis on \(\K(E)\) gives property \((M)\) of \(E\)
and the compact approximation net required by \cite[Theorem 4.7]{DL}.
By \cite[Proposition 3.10]{D}, \(E\) has scalar \(n\)-polynomial
property \((M)\). Lemma \ref{lem:lift} upgrades this to vector
\(n\)-polynomial property \((M)\) for \((E,F)\), while Lemma
\ref{lem:cd} gives \(n=cd(E,F)\). Applying \cite[Theorem 4.7]{DL}
proves the conclusion.
\end{proof}

\section{Deep upgrade attempt and limitation}

The finite-dimensional proof suggests choosing norming functionals
\(f_\alpha\) for an arbitrary property-\((M)\) range \(F\). Banach--Alaoglu
only gives a weak-star convergent subnet. It does not imply
\[
 (f_\alpha-f)(P(x_\alpha))\longrightarrow0,
\]
because both factors move with \(\alpha\), and at the critical degree
\(P(x_\alpha)\) need not be weakly null. This is exactly the point at which
the strict inequality \(n<cd(E)\) was used in Proposition 4.8 of \cite{DL}.

I also tested compact-operator approximation on \(F\). A fixed compact
operator makes the norming-functional argument finite-dimensional on its
image, but no such operator uniformly approximates the range of an arbitrary
critical-degree polynomial. The stronger \(M_\infty\)-decomposition estimate
does handle important infinite-dimensional targets, but that is already the
mechanism of Propositions 2.7--2.8 in \cite{DL}. No general property-\((M)\)
upgrade is claimed.

\section{Verification and novelty status}

The proof uses no computation. The separate \texttt{verification.md} audits
critical-degree equality, complex Hahn--Banach scalarization, limsup subnets,
and every imported theorem.

Local indexes were searched for arXiv:1102.3850, the exact title,
\(n\)-polynomial property \((M)\), the equality \(cd(E,F)=cd(E)\), and
finite-dimensional ranges. Bounded web searches used the same exact phrases
and checked the source, the scalar predecessor arXiv:1005.1260, and the later
related arXiv:1512.08741. No exact finite-dimensional endpoint statement was
found. This is bounded evidence, not a guarantee of novelty.

\section*{Human-review recommendation}

Promote as a likely-valid partial result after checking the two-line subnet
scalarization in Lemma \ref{lem:lift} and the cited theorem numbers. The main
claim should remain explicitly limited to finite-dimensional \(F\).

\begin{thebibliography}{9}
\bibitem{D}
V. Dimant,
\emph{\(M\)-ideals of homogeneous polynomials},
Studia Math. 202 (2011), 81--104;
arXiv:1005.1260.

\bibitem{DL}
V. Dimant and S. Lassalle,
\emph{\(M\)-structures in vector-valued polynomial spaces},
J. Convex Anal. 19 (2012), 685--711;
arXiv:1102.3850.
\end{thebibliography}

\end{document}
